\begin{afterword}
\summaryhead

The post begins with psychophysics. People can rate how loud a sound is on an unbounded scale, and when given a standard sound with a fixed number, the modulus, their ratings agree. Without a modulus each person picks a private scale: the ratios between sounds agree, but the raw numbers vary enormously.

Kahneman, Schkade and Sunstein found the same in punitive damages. Their 867 jury-eligible respondents agreed on how outrageous a defendant's conduct was and how much to punish it, on bounded scales, and agreed on the ranking of dollar awards. The dollar amounts themselves were close to unpredictable, even as medians of twelve. So an award is an expression of outrage on a scale with no modulus.

The author then proposes that forecasts of when human-level AI will arrive are the same. A researcher who once said ``Five hundred years'' was, the author guesses, expressing how hard the problem feels, or how the forecaster feels about AI, and adding ``years.''

\responsehead

The first two thirds of the post report the jury study accurately. The figures match the paper, and the analogy with magnitude scaling is the paper's own. Two claims go beyond it. The dollar award is not ``completely unpredictable'': the case explained 6 per cent of the variance in raw awards, and 42 per cent in their logarithms, close to the 51 per cent for ranks. The authors agree that log awards are ``reasonably predictable,'' while arguing that parties pay in dollars. And the step from survey respondents to real juries is the authors' own, taken with more caution in one version of their work than in the other: their juries were medians of individual answers to short vignettes, with no deliberation.

The last third makes a new claim with none of the jury study's support. The attitude-expression reading of punitive damages rested on a pattern: people who disagree wildly about dollars agree about outrage and about the order of cases. The post shows no such pattern for AI forecasts. It offers the spread of answers the author has seen and one remembered conversation. Spread alone cannot separate the two readings, because honest forecasters facing a date the post itself calls ``not very predictable'' would also disagree widely. The guess about what ``Five hundred years'' meant is offered as a guess, and should be read as one.

The thesis may still be right. A 2016 survey of machine-learning researchers found answers that shifted with the wording of the question, including a 50 per cent chance at 45 years for one question and at 122 years for another that its authors call logically similar. That is the kind of evidence the post needed. It came later, and the post does not have it.

In short: an accurate account of the jury study, with its unpredictability somewhat overstated, followed by a claim about AI forecasts that rests on one anecdote and on a spread of answers that uncertainty alone could explain.
\end{afterword}
